\chapter{Conclusions}

The Conclusion chapter synthesizes the dissertation, reinforcing its key findings and contributions while avoiding repetition of previous sections’ phrasing.

\section{Content Elements}
\begin{itemize}
    \item Restate the research objectives and primary results, emphasizing their significance.
    \item Summarize the study’s contributions to academia or practice.
    \item Offer specific recommendations for future research or applications.

\end{itemize}

\section{Writing Tips}
\begin{itemize}
    \item Use concise, impactful language to highlight the study’s value.
    \item Echo the research questions from the Introduction for a cohesive narrative.
    \item Avoid introducing new data or unaddressed perspectives.
    \item Consider outlining the study’s long-term impact, such as potential influences on policy, industry, or academic development.

\end{itemize}
